\section{はじめに}

日中に眠気を感じる人は多い．厚生労働省の令和元年国民・栄養調査報告\cite{[1]}によると，男性の32.3\%，女性の36.9\%が日中に眠気を感じていることが報告されている．適度な仮眠は集中力を高めたり，脳を活性化させることが知られているが，業務の形態や進捗によっては仮眠が許されないことも多いし，意図しない居眠りが業務の遅れや事故に繋がることもある．特に昼食後のデスクワークにおける居眠りを防ぐためには，システム面から居眠りの検知・防止ができることが望ましい．

これまで，居眠り（覚醒度低下）検知システムの開発は，運転中における居眠りなどを対象として数多く行われてきた．既存の方法は，体動の変化から居眠りを検知する方法と生理指標から居眠りを検知する方法の大きく二つに分けられる．前者には，画像処理による人の瞬目の変化からの居眠り検知\cite{[2]}，加速度センサによる頭の動きからの居眠り検知\cite{[3]}，座席の背もたれに設置した導波管型センサからの居眠り検知\cite{[4]}などがある．後者には，心拍の時系列解析による居眠り検知\cite{[5]}，脳波からの居眠り検知\cite{[6]}などがある．
いずれも眠気と相関関係のある物理量を計測し，居眠りを推測している．特に脈拍センサによる居眠りの検知方法は，ストレス値を定量的に評価したり，視覚化したりできるため，先行研究も多い\cite{[5]}\cite{[7]}\cite{[9]}．ところで，このセンサは，指に正しく装着されていない場合や意図せず指から外れていた場合に乱数値が出力されることがある．しかし，そのような場合にそれが使用者の誤った使用方法に起因していることをセンサ側から検知できないという技術的課題がある．

そこで本プロジェクトでは，体動の変化からの居眠り検知の新たな方法として，光センサを用いた居眠り検知・防止を提案する．これは，昼食後のデスクワークにおける居眠り検知・防止を想定している．具体的には，頭の動きを光センサで捉え，頭の影と環境光との光量の差から居眠りを検知する．居眠りを検知すると，モータによって振動を起こし，使用者を起床させる．また，居眠り検知の既存の方法の技術的課題の解決として，使用法に誤りがあることを検出し，使用者に警告する機能を加える．

また，本プロジェクトでは，脈拍センサを用いて使用者のストレス状態を視覚化した上で，光センサによる居眠り検知が正しく行われているかを検証した．具体的には，睡眠状態の解析に用いられるローレンツプロット（LP: Lorenz Plot）\cite{[7]}\cite{[8]}をProcessingによって表示し，使用者のストレス状態を幾何学的に視覚化した．さらに，松本佳昭らの論文で提案されている評価指標F-stress\cite{[9]}，心拍間隔（RRI: R-R Interval），二つの光センサの変動について使用者の状態ごとに記録し，比較して解析した．LP情報と実験の結果より，本装置によって使用者の居眠り検知・防止ができることを明らかにする．さらに，本装置はF-stressでは取得できない使用者のうとうと状態をグラフ上に波形として捉えられることが可能であることも明らかにする．本システムは，この点において既存の評価指標よりも有効性があると言える．また，誤使用時や暗所時においても想定通りの動作が正しく行われていることを明らかにする．